% Пример статьи для журнала КИО
\documentclass[intlimits,twoside,a4paper,11pt]{article}


% текст статьи рекомендуется писать в кодировке utf-8
\usepackage[utf8]{inputenc}

\usepackage{amsmath,amssymb,amsthm}


% Подключение пакета с шаблоном журнала.
% Необзятальная опция eqsecnum указывает, что в каждом разделе формулы нумеруются заново.
% Опция authorsinfo заставляет выводиться информацию об авторе
\usepackage[eqsecnum,authorsinfo]{kioj4}
% Здесь можно подключить свои пакеты. Некоторые пакеты работают только если их подключить перед предыдущей строкой \usepackage[eqsecnum]{kioj4}
% \usepackage[•]{•}
\usepackage{array}


% Указываем номер страницы, с которой начинается статья, можно не указывать, будет нумероваться с первой.
\setcounter{page}{41}

%рубрика в журнале, обычно не известна при подаче статьи, поэтому не указываем
\journalsectionnoimage{informatics}
%\journalsection{informatics}
%\journalsection{information-systems}
%\journalsection{software-engineering}
%\journalsection{computer-in-education}
%\journalsection{popular-science-articles}
%\journalsection{programming-practice}
%\journalsection{editorial-column}
%\journalsection{specialists-training-new-teaching-methods}
%\journalsection{specialists-training-training-programs}
%\journalsection{specialists-training-professional-standards}
%\journalsection{open-problems-for-young-scientists}
%\journalsection{empty}

%указывам год, номер выпуска неизвестен, поэтмоу ставим прочерк, указываем edu, что означает журнал "компьютерные инструменты в образовании".
\issue{2026}{-}{edu}
% Указываем УДК для статьи
\udknumber{511.331.1}
% Указываем DOI для статьи, можно не указывать
\doinumber{10.1000/182}

%Указываем название статьи в обычном регистре в квадратных скобках, потом заглавными буквами в фигурных скобках. Второй вариант будет использоваться только в заголовке, первый вариант будет использоваться во всех других местах, в частности, при цитировании.
\title[Исследование коэффициентов конечных рядов Дирихле, обращающихся в ноль в некоторых нулях дзета-функции Римана]{ИССЛЕДОВАНИЕ КОЭФФИЦИЕНТОВ КОНЕЧНЫХ РЯДОВ ДИРИХЛЕ, ОБРАЩАЮЩИХСЯ В НОЛЬ В НЕКОТОРЫХ НУЛЯХ ДЗЕТА-ФУНКЦИИ РИМАНА}

%Информация об авторах. Каждый автор указывается отдельно с помощью трех инструкций \authorI, \authorIaddr, \authorIemail. Номера авторов могут быть I, II, III, IV, V.
%Один из авторов должен быть основным получателем корреспонденции, в этом случае вместо команды \author*email нужно использовать команду \author*emailEnvelope
%в команде \author* указывается сначала имя автора с инициалами в квадратных скобках, потом указывается полное имя автора. Полное имя используется только при выводе расширенной информации об авторе.
%в команде \affil указывается через запятую список организаций, которые представляет автор. Сами организации указываются ниже командой \affiliation.
%\authorIaddr и \affilitaion различаются тем, что первое является кратким расском об авторе, в отличие от второго, где указывается только организация.
\authorI[Черепанов Р.~П.\affil{1}]{Черепанов Роман Павлович}
\authorIpos{Магистрант}
\authorIemailEnvelope{romacherepanov2002@gmail.com}

\authorII[Субботин М.~О.\affil{1}]{Субботин Максим Олегович}
\authorIIpos{Аспирант}
\authorIIemail{mosubbotin@etu.ru}

\authorIII[Емельянов Д.~С.\affil{2}]{Емельянов Дмитрий Сергеевич}
\authorIIIpos{Магистрант}
\authorIIIemail{dima31120251@gmail.com}

\affiliation{1}{СПбГЭТУ, Санкт-Петербургский государственный электротехнический университет «ЛЭТИ» им. В. И. Ульянова (Ленина), ул. Профессора Попова, 5, корп. 3, 197022, Санкт-Петербург, Россия}
\affiliation{2}{ИТМО, Университет ИТМО, Кронверкский пр., 49, Санкт-Петербург, 197101, Россия}

\begin{document}
    \maketitle

    \begin{abstract}
        В 2013 году Ю. В. Матиясевич численно исследовал конечные ряды Дирихле, обращающиеся в 
        ноль в нескольких нетривиальных нулях дзета-функции Римана. Он зафиксировал первый коэффициент 
        таких рядов равным 1 и обнаружил, что начальные коэффициенты рассматриваемых рядов Дирихле весьма 
        близки к коэффициентам знакопеременной дзета-функции (эта-функции). 


        Это исследование в настоящей работе расширено путём фиксации 
        нескольких коэффициентов таких конечных рядов Дирихле, а также рассмотрением рядов, построенных по нулям $L$-функций.
        Выявлены некоторые закономерности в поведении оставшихся начальных коэффициентов. 
        В частности, обнаружено, что для ряда, построенного по нулям дзета-функции, коэффициенты приближаются
        коэффициентами произведения $\zeta(s)\sum_{k=1}^{f+1}{b_k k^{-s}}$, где $f$ --- число фиксированных коэффициентов,
        а значения $b_k$ определяются из коэффициентов самого ряда. 


        Результаты данного анализа дают новые представления о связи между дзета-функцией Римана и теорией чисел.




        \keywords дзета-функция, линейная алгебра, аналитическая теория чисел.
        %Команда \autocitationexample автоматически генерирует пример цитирования статьи, либо можно самостоятельно написать цитирование с помощью команды \citationexample
        \autocitationexample
        %\citationexample Андреев А.~А., Борисов Б.~Б. Шаблон статьи для журнала КИО. Нигде не опубликовано
        \acknowledgements Авторы выражают благодарность Юрию Владимировичу Матиясевичу 
        за постановку задачи и консультации по ходу исследования. 
    \end{abstract}

    \section{ВВЕДЕНИЕ}

        В 1735 году Леонард Эйлер опубликовал решение базельской задачи: $$\sum_{n=1}^\infty \frac{1}{n^2} = \frac{\pi^2}{6}.$$
        Затем он продолжил изучать свойства функции вещественной переменной $f(x) = \sum_{n=1}^\infty n^{-x}$
        и в 1737 году обнаружил тождество $$\sum_{n=1}^\infty n^{-x} = \prod_{p \in \mathbb{P}}\frac{1}{1-p^{-x}}.$$
        Это тождество известно как тождество Эйлера.

        Бернхард Риман в 1859 году ввел дзета-функцию $$\zeta(s) = \sum_{n=1}^\infty n^{-s},$$ обобщив функцию,
        которую рассматривал Эйлер, на комплексный аргумент. Данный ряд сходится только для $\mathrm{Re}(s) > 1$.
        Для всех остальных значений аргумента построено аналитическое продолжение, которое имеет полюс только
        в точке $s = 1$. В этой же публикации 1859 года Риман выдвинул гипотезу, согласно которой
        все нетривиальные нули дзета-функции имеют вещественную часть, равную $1/2$. Гипотеза Римана
        является одной из ключевых задач аналитической теории чисел. Связь между гипотезой Римана и теорией чисел
        демонстрирует тождество Эйлера, и в предположении о верности гипотезы доказывается ряд теорем
        о распределении простых чисел~\cite{sviridov, matveenko} и росте арифметических функций~\cite{akim}.
       
        Несмотря на значительный прогресс в понимании асимптотического поведения $\zeta(s)$, 
        вопросы о точных приближениях этой функции в различных областях комплексной плоскости остаются актуальными. 
        В частности, конструкция конечных рядов Дирихле, аппроксимирующих $\zeta(s)$ и точно воспроизводящих 
        её поведение в окрестностях первых нетривиальных нулей, открывает возможности как для численной проверки 
        гипотезы Римана, так и для развития методов в теории $L$-функций. Здесь и далее для нумерации нулей используются их модули. 
        Первыми нетривиальными нулями дзета-функции называются нетривиальные нули с наименьшей по модулю мнимой частью. 
        При равенстве модулей первым нулем считается нуль с положительной мнимой частью.

        Перспективный подход к построению таких приближений был предложен в работах Ю. В. Матиясевича~\cite{Chernov2015, de430}. 
        Суть метода заключается в построении конечного ряда Дирихле 
        \begin{equation}
            R(s) = \sum_{k=1}^{N} a_{k} k^{-s},
        \end{equation}
        коэффициенты которого подбираются таким образом, чтобы ряд $R(s)$ обращался в ноль в первых $N$ нетривиальных нулях дзета-функции. 
        Данная задача сводится к решению системы линейных уравнений, где каждое уравнение соответствует условию 
        $R(s_j) = 0$ для $j$-го нуля $s_j$. Для обеспечения нетривиальности решения дополнительно фиксируется значение одного из коэффициентов 
        (например, $a_1=1$). Такая система имеет матричную запись $\mathbf{S} \mathbf{a} = \mathbf{b}$, 
        где $\mathbf{S}$ — квадратная матрица размера $N \times N$, $\mathbf{a}$ — вектор неизвестных коэффициентов, 
        а $\mathbf{b}$ — вектор свободных членов. 

        Настоящая работа обобщает указанный подход на случай фиксации нескольких коэффициентов ряда $R(s)$.
        Формально, если требуется обратить $R(s)$ в ноль в $n$ точках и зафиксировать $f$ коэффициентов, 
        то общее число параметров определяется соотношением $N = n + f$. Соответствующая система линейных уравнений 
        при этом содержит $f$ строк, задающих фиксированные значения коэффициентов, и $n$ строк, отвечающих за 
        обращение $R(s)$ в ноль в заданных точках. Выбор фиксируемых коэффициентов влияет на 
        обусловленность системы и, следовательно, на численную устойчивость решения. 

        Для иллюстрации рассмотрим случай трёх фиксированных коэффициентов $a_1=1$, $a_2=0$, $a_3=0.5$. Тогда система уравнений принимает вид:
        \begin{equation}\label{eq:second_system}
        \begin{bmatrix}
        1 & 0 & 0 & 0 & \dots & 0 \\
        0 & 1 & 0 & 0 & \dots & 0 \\
        0 & 0 & 1 & 0 & \dots & 0 \\
        1^{-s_1} & 2^{-s_1} & 3^{-s_1} & 4^{-s_1} & \dots & N^{-s_1}\\
        1^{-s_2} & 2^{-s_2} & 3^{-s_2} & 4^{-s_2} & \dots & N^{-s_2}\\
        \vdots & \vdots & \vdots & \vdots & \ddots & \vdots \\
        1^{-s_{n}} & 2^{-s_{n}} & 3^{-s_{n}} & 4^{-s_{n}} & \dots & N^{-s_{n}} \\
        \end{bmatrix}
        \begin{bmatrix} 
        a_1 \\ a_2 \\ a_3 \\ a_4 \\ \vdots \\ a_N 
        \end{bmatrix}
        =
        \begin{bmatrix} 
        1 \\ 0 \\ 0.5 \\ 0 \\ \vdots \\ 0 
        \end{bmatrix},
        \end{equation}
        где $N = n + 3$. Данная конструкция естественным образом обобщается на произвольное число $f$ фиксированных коэффициентов.
  
        В вычислительных экспериментах работы используются значения нетривиальных нулей дзета-функции Римана доступные в базе данных Университета Дикина~\cite{pavlov}. 
        Анализ полученных результатов позволяет выявить новые закономерности в 
        распределении коэффициентов $a_k$, которые могут иметь значение для дальнейших исследований в области аналитической теории чисел.


    \section{ИССЛЕДОВАНИЕ ТОЧНОСТИ ВЫЧИСЛЕНИЙ В ЗАВИСИМОСТИ ОТ ТОЧНОСТИ НЕТРИВИАЛЬНЫХ НУЛЕЙ}
        Числа в компьютерной алгебре представляются неточно, поэтому для работы с высокоточными нетривиальными нулями 
        использовалась интервальная арифметика, реализованная в библиотеке Arb~\cite{iers2010} пакета FLINT~\cite{standish}. 
        Нетривиальный нуль $x$ заключен в интервал $l \leqslant x \leqslant u$. 
        Был построен ряд Дирихле с использованием 100 нулей дзета-функции с точностью 200 знаков. 
        На рисунке~\ref{fig-oylog} представлен график длин интервалов $u-l$ и $\log(u-l)$ для каждого коэффициента ряда. 
        Для конечного ряда с $n$ коэффициентами при коэффициентах с номерами $[0.7n, 0.85n]$ наблюдается падение точности. 
        Такое поведение соотносится с тем, что до номера $0.7n$ коэффициенты принимают чередующиеся значения ${-1, 1}$, после чего значения коэффициентов убывают по модулю и приближаются к 0 после $0.85n$. 
        Также на рис.~\ref{fig-log50} и рис.~\ref{fig-log150} представлены результаты для рядов Дирихле с 50 и 150 нулями дзета-функции. 


        

            Рассмотрим минимальные, средние и максимальные длины интервалов $u-l$ для разных точностей входных данных. 
            Количество «верных» цифр, обеспечиваемое длинами интервалов, можно оценить как $-\lfloor\log_{10}(u - l)\rfloor - 1$. 
            На рисунке~\ref{fig-interval1} наблюдается линейная зависимость точности нулей и длин интервалов коэффициентов, 
            поэтому для увеличения количества «верных» цифр коэффициентов ряда Дирихле на $N$ достаточно увеличить точность нулей дзета-функции на $N$ «верных» цифр. 


        


    \section{ИССЛЕДОВАНИЕ КОЭФФИЦИЕНТОВ РЯДОВ ДИРИХЛЕ С ИСПОЛЬЗОВАНИЕМ НУЛЕЙ ДЗЕТА-ФУНКЦИИ}  
        Основное внимание в данном разделе сосредоточено на начальных коэффициентах ряда Дирихле, 
        поскольку к середине ряда и последующим членам наблюдается монотонный рост их модуля, 
        что затрудняет прямое исследование полной последовательности коэффициентов.
        Все ряды в этом разделе построены с использованием 400 нулей дзета-функции с точностью 600 цифр. 

        \subsection{Классы делимости}

        Первые 80 коэффициентов ряда с 2 и 3 фиксированными начальными коэффициентами представлены на рис.~\ref{fig-c1} ($a_1=1, a_2=0$) 
        и рис.~\ref{fig-c2} ($a_1=1, a_2=0, a_3=0.5$). На рисунках горизонтальными пунктирными линиями помечены значения начальных коэффициентов
        с близкими значениями. Таким образом на рис.~\ref{fig-c1} получено 4 таких линий, а на рис.~\ref{fig-c2} их 6. 
        Была найдена следующая закономерность: коэффициенты лежат на одной линии, когда множества их делителей среди
        чисел $(2, 3, \dots, f+1)$ совпадают. Каждому такому допустимому множеству делителей соответствует \textit{класс делимости}.

        С увеличением количества фиксированных начальных коэффициентов $f$ максимальное количество классов делимости $c_f$ увеличивается или остается неизменным.
        В процессе проведения экспериментов были выявлены следующие свойства:
        \begin{enumerate}
            \item Когда $f+1$ является степенью простого числа $p^k$, только классы, содержащие $p^{k-1}$, делятся на два класса.
            \item Когда $f+1$ не является степенью простого числа, количество классов не увеличивается. 
            Например, при $f+1=6$, делимость на 6 определяется делимостью на 2 и 3.
        \end{enumerate}
        В результате для $f$ фиксированных коэффициентов количество классов $c_f$ определяется по формуле:

    \begin{equation}
        c_f=
        \begin{cases}
            2, & \mbox{если } f=1; \\
            \frac{k+1}{k}\cdot c_{f-1}, & \mbox{если }f+1\mbox{ --- простое число в степени }k; \\
            c_{f-1}, & \mbox{если }f+1\mbox{ --- составное число}.\\
        \end{cases}
    \end{equation}

    В таблице~\ref{table-f} представлены значения количества классов делимости $c_f$ для первых значений $f$. 

    
        

    \begin{table}[H]
        \centering
        \caption{Количество классов делимости $c_f$.}
        \label{table-f}
        \small
        \begin{tabular}{|p{5mm}|c|}
            \hline
            \centering
            $f$ & $c_f$ \\
            \hline
            \centering
            1 & 2 \\
            \hline
            \centering
            2 & 4 \\
            \hline
            \centering
            3 & 6 \\
            \hline
            \centering
            4 & 12 \\
            \hline
            \centering
            5 & 12 \\
            \hline
            \centering
            6 & 24 \\
            \hline
            \centering
            7 & 32 \\
            \hline
            \centering
            8 & 48 \\
            \hline
            \centering
            9 & 48 \\
            \hline
        \end{tabular}
    \end{table}

    Для других наборов фиксированных индексов коэффициентов поведение построенных рядов может различаться. 
    Если индексы из одного и того же класса делимости ограничены разными значениями, то величины значений 
    ряда оказываются больше ожидаемых, как показано на рисунке~\ref{fig-c3}.

    

    Другим возможным вариантом поведения результирующего ряда является уменьшение количества 
    различных по величине значений коэффициентов. При определенных ограничениях результирующий ряд 
    оказывается подобным ряду, построенному на основе меньшего числа ограничений, как показано на рисунке~\ref{fig-c4}, 
    где 3 фиксированных коэффициента приводят к 4 классам величин коэффициентов. 


    В некоторых случаях решение расходится, даже если наложенные ограничения относятся к разным классам делимости (\ref{fig-c5})

   


    \subsection{Вычисление значений многочленов, описывающих полученные ряды}

    Для одного фиксированного коэффициента начальные коэффициенты результирующего ряда близки к результату произведения:
    \begin{equation}
        (b_1+b_2\cdot2^{-s})\zeta(s)
    \end{equation}

    Обобщением этого утверждения является гипотеза, согласно которой начальные коэффициенты результирующего ряда при $f$ фиксированных коэффициентах напоминают следующее произведение:
    \begin{equation}\label{eq:ff}
        \zeta(s)\sum_{k=1}^{f+1}{b_k k^{-s}}
    \end{equation}

    Например, при фиксированных коэффициентах $a_1=1$ и $a_2=0$ (см. рис.~\ref{fig-c1}), первые 6 коэффициентов ряда Дирихле 
    $a_1=1$, $a_2=0$, $a_3=-0.5$, $a_4=0$, $a_5=1$ и $a_6=-1.5$ образуют следующие слагаемые:

    \begin{equation}
        1^{-s} + 0 \cdot 2^{-s} + (-0.5) \cdot 3^{-s} + 0 \cdot 4^{-s} + 1 \cdot 5^{-s} + (-1.5) \cdot 6^{-s}
    \end{equation}

    Первые 6 слагаемых произведения \ref{eq:ff} (при $f=2$) будут следующими:
    \begin{equation}
        b_1 \cdot 1^{-s} + (b_1 + b_2) \cdot 2^{-s} + (b_1 + b_3) \cdot 3^{-s} + (b_1 + b_2) \cdot 4^{-s} + b_1 \cdot 5^{-s} + (b_1 + b_2 + b_3) \cdot 6^{-s}
    \end{equation}

    Соответствующая система уравнений имеет вид:

    \begin{equation*}
        \begin{cases}
            1 = b_1, \\
            0 = b_1 + b_2,\\
            -\frac{1}{2} = b_1 + b_3,\\
        \end{cases} \implies 
        \begin{cases}
            b_1 = 1, \\
            b_2 = -1, \\
            b_3 = -\frac{3}{2}.
        \end{cases}
    \end{equation*}

    Таким образом, ряд $R(s)$ схож с $(1 - 2^{-s} - 1.5 \cdot 3^{-s}) \zeta(s)$ на начальных коэффициентах.  

    В общем случае коэффициент результирующего ряда $a_i$ равен сумме коэффициентов $b_j$, где $j | i$.  
    Такая система уравнений позволяет вычислить все значения $b_j$. Например, для 5 фиксированных коэффициентов 
    система уравнений для 6 коэффициентов $b_k$ имеет вид:

    \begin{equation*}
        \begin{cases}
            a_1 = b_1  \\
            a_2 = b_1 + b_2 \\
            a_3 = b_1 + b_3 \\
            a_4 = b_1 + b_2 + b_4 \\
            a_5 = b_1 + b_5 \\
            a_6 = b_1 + b_2 + b_3 + b_6 \\
        \end{cases}
    \implies
    \begin{cases}
            b_1 = a_1 \\
            b_2 = a_2 - b_1 \\
            b_3 = a_3 - b_1 \\
            b_4 = a_4 - b_1 - b_2 \\
            b_5 = a_5 - b_1 \\
            b_6 = a_6 - b_1 - b_2 - b_3 \\
        \end{cases}
    \implies
    \begin{cases}
            b_1 = a_1 \\
            b_2 = a_2 - a_1 \\
            b_3 = a_3 - a_1 \\
            b_4 = a_4 - a_2  \\
            b_5 = a_5 - a_1 \\
            b_6 = a_6 - a_2 - a_3 + a_1
        \end{cases}
    \end{equation*}

    После проведенного анализа системы уравнений для $b_k$ можно объяснить неожиданное расхождение (рис.~\ref{fig-c5}). 
    При 4 фиксированных коэффициентах имеется 5 коэффициентов $b_k$. Система уравнений имеет вид:

    \begin{equation}
        \begin{cases}
            a_{10} = b_1 + b_2 + b_5; \\
            a_{20} = b_1 + b_2 + b_4 + b_5; \\
            a_{30} = b_1 + b_2 + b_3 + b_5; \\
            a_{60} = b_1 + b_2 + b_3 + b_4 + b_5.
        \end{cases}
    \end{equation}

    Можно заметить, что ранг матрицы равен 3:

    \begin{equation}
        \begin{bmatrix}
            1 & 1 & 0 & 0 & 1 \\
            1 & 1 & 0 & 1 & 1 \\
            1 & 1 & 1 & 0 & 1 \\
            1 & 1 & 1 & 1 & 1
        \end{bmatrix}
    \end{equation}

    Следовательно, на такую систему нельзя произвольно наложить 4 ограничения.

    Следует также отметить, что фиксация коэффициентов из одного и того же класса дополнительно 
    уменьшает ранг матрицы, что приводит к несовместности системы, если ранг расширенной 
    матрицы не уменьшается соответствующим образом.


    \section{$L$-функции}

    В данном разделе рассматриваются функции $L_{3,2}$ и $L_{5,4}$, ряды Дирихле которых выглядят как ряды $\zeta$-функции,
    умноженные на характеры Дирихле $\chi_3(2,\cdot)$ и $\chi_5(4,\cdot)$ соответственно (эти характеры определяются равенствами 
    $\chi_3(2,2)=-1$, $\chi_5(4,2)=-1$, $\chi_5(4,3)=-1$ и $\chi_5(4,4)=1$).
    
    Проведено исследование по построению рядов с 
    использованием нулей этих L-функций~\cite{bizouard11}. В следующих вычислениях использовались 300 пар нулей с точностью 600 цифр. 

    \subsection{Примеры фиксирования коэффициентов, $L_{3,2}$}
    
    При фиксировании первого коэффициента получается следующее решение (рис.~\ref{fig-l-1}), 
    где период значений ряда равен 3, и значения коэффициентов ряда совпадают с характером Дирихле $\chi_2\ \mathrm{mod}\ 3$ 
    при подстановке в него соответствующих индексов.

    Решение при зафиксированных первых двух коэффициентах ($a_1=1$ и $a_2=3$) представлено на рисунке~\ref{fig-l-2}, период ряда будет равен 6.

    
 

    Была выдвинута следующая гипотеза: первые члены ряда, построенного по нулям $L$-функции (с индексами 3, 2) с первыми
    $k$ фиксированными коэффициентами, отвечают формуле

    \begin{equation}
        F(s) = \sum_{i=1}^{}{a_{i\ \mathrm{mod}\ M}\cdot i^{-s}},
    \end{equation}

    где $M$ --- длина периода и $a_i$:

    \begin{equation}
        a_i=\sum_{j|i}^{f}{b_j \cdot \chi(i/j)},
    \end{equation}

    где $\chi$ --- характер Дирихле, определяющий $L$-функцию, $f$ --- количество 
    зафиксированных коэффициентов ряда. При этом период равен $\mathrm{НОК}(1,2,\dots,f)\cdot N$, где $N$ --- период характера. 

    Отсюда для первых 6 коэффициентов в примере выше (рис.~\ref{fig-l-2}) система уравнений принимает вид:

    \begin{equation}
        \begin{cases}
            a_{1} = b_1 \\
            a_{2} = -b_1 + b_2 \\
            a_{3} = 0 \\
            a_{4} = b_1 - b_2 \\
            a_{5} = -b_1 \\
            a_{6} = 0
        \end{cases}
    \end{equation}

    \subsection{Примеры фиксирования коэффициентов, $L_{5, 4}$}

    По аналогии с предыдущим случаем, при фиксировании первого коэффициента график совпадает с характером
    Дирихле $\chi_4\ \mathrm{mod}\ 5$  (рис.~\ref{fig-l-3}). При фиксировании 4 коэффициентов (рис.~\ref{fig-l-4}) длина периода составила
    НОК(2, 3, 4) $\cdot$ 5 = 60.

     
 

    Уравнения для первых 6 коэффициентов:

    \begin{equation}
        \begin{cases}
            a_{1} = b_1 \\
            a_{2} = -b_1 + b_2 \\
            a_{3} = -b_1 + b_3 \\
            a_{4} = b_1 - b_2 + b_4 \\
            a_{5} = 0 \\
            a_{6} = b_1 - b_2 - b_3
        \end{cases}
    \end{equation}
    
    \section{ЗАКЛЮЧЕНИЕ}

    Была выявлена зависимость длины интервалов, в которые заключены коэффициенты результирующего ряда, от точности
    используемых для его построения нулей дзета-функции. При увеличении количества верных цифр нулей в $n$ раз одновременно
    росло количество точно вычисленных цифр коэффициентов результирующего ряда в $n$ раз.

    В ходе исследования были выявлены несколько конфигураций индексов, влияющих на ряды, построенные с использованием 
    начальных нулей дзета-функции. При наложении $k$ ограничений на значения коэффициентов ряда начальные коэффициенты
    результирующего ряда напоминают результат произведения
    \begin{equation}
        \zeta(s)\left(b_1\cdot 1^{-s} + \dots + b_{k+1} \cdot (k+1)^{-s}\right) .
    \end{equation}

    Если система уравнений, образованная наложенными ограничениями, несовместна, решение расходится.

    В ходе исследования было установлено, что ряды, построенные по нескольким нулям функций $L_{5, 4}$ и $L_{3, 2}$, обладают
    похожим свойством. При фиксации $k$ коэффициентов наблюдается периодичность, и для коэффициентов в периодах выполнено равенство
    \begin{equation}
        a_i = \sum_{j|i}^{f}{b_j \cdot \chi(i/j)}.
    \end{equation}

    %Оформляем литературу, 9 означает, что у нас не более 9 пунктов в списке, т.е. все номера состоят из не более, чем одной цифры. Если пунктов много, т.е. номера могут быть двузначными, пишем 99.
    \bibliographystyle{ugost2008} 
    \bibliography{references}




    %указываем, что теперь начинается английская часть

    \begin{translatedpart}
        %указываем название и авторов на английском
        \title{Study of coefficients of finite Dirichlet series vanishing at some zeros of Riemann's zeta function}
        \authorI[Cherepanov R.~P.\affil{1}]{Cherepanov Roman Pavlovich}
        \authorIpos{Master's degree student}
        \authorIemailEnvelope{romacherepanov2002@gmail.com}

        \authorII[Subbotin M.~O.\affil{1}]{Subbotin Maksim Olegovich}
        \authorIIpos{PhD student}
        \authorIIemail{mosubbotin@etu.ru}  


        \authorIII[Emelyanov D.~S.\affil{2}]{Emelyanov Dmitry Sergeevich}
        \authorIIIpos{Master's degree student}
        \authorIIIemail{dima31120251@gmail.com}

        \affiliation{1}{Saint Petersburg Electrotechnical University, 5, building 3, Professora Popova st., 197022, Saint Petersburg, Russia}
        \affiliation{2}{ITMO University, Kronverksky Ave., 49, Saint Petersburg, 197101, Russia}

        %даем команду на печать английского заголовка
        \maketranslatedtitle

        \begin{abstract}
            In 2013, Yu.~V.~Matiyasevich numerically investigated finite Dirichlet series vanishing at 
            several non-trivial zeros of the Riemann zeta function. 
            He fixed the first coefficient of such series to~1 and found 
            that the initial coefficients of the Dirichlet series under 
            consideration are very close to the coefficients of the 
            alternating zeta function (the eta function).


            This investigation is extended in the present work by fixing several coefficients 
            of such finite Dirichlet series, as well as by considering 
            series constructed from the zeros of $L$-functions. 
            Certain patterns in the behavior of the remaining initial coefficients 
            have been identified. In particular, it is found that for the 
            series constructed from the zeros of the zeta function, the 
            coefficients approximate those of the product 
            $\zeta(s)\sum_{k=1}^{f+1}{b_k k^{-s}}$, where~$f$ is the number of 
            fixed coefficients, and the values~$b_k$ are determined from the coefficients of the series itself.


            The results of this analysis provide new insights into the connection between the Riemann zeta function and number theory.

            \keywords zeta function, linear algebra, analytic number theory.
            \autocitationexample
            %\citationexample ???
            \acknowledgements The authors express their gratitude to Yuri~V.~Matiyasevich for posing the problem and for consultations during the research.
        \end{abstract}

        

        \additionalinfo{Received October 7, 2001, The final version: December 28, 2015}

    \end{translatedpart}

    

\end{document}
